% Solutions of 17-balanced-trees.
\section*{Chapter \ref{ch:b1:balanced-trees} --- Binary Search Trees and Balanced Trees}

\begin{solution}{exo:b1:balanced-trees:1}
The tree has root 50, children 30 and 70; 30 has children 20 and 40, and 35 is the left child of 40; 70 has children 60
and 80. In order: 20, 30, 35, 40, 50, 60, 70, 80. \omterm{def:b1:balanced-trees:traversal}{Preorder}: 50, 30, 20, 40, 35, 70, 60, 80. \omterm{def:b1:balanced-trees:traversal}{Postorder}: 20, 35, 40, 30,
60, 80, 70, 50.
\end{solution}

\begin{solution}{exo:b1:balanced-trees:2}
Root 8; left 3 with children 1 and 6, and 6 with children 4 and 7; right 10 with right child 14, whose left child is 13.
The height is 3 (paths 8, 3, 6, 4 and 8, 10, 14, 13). Deleting 3, which has two children: its \omterm{def:b1:balanced-trees:ordered}{successor} 4, the minimum
of its right \omterm{def:b1:induction-and-termination:rooted-tree}{subtree}, takes its place, and 6 keeps only its right child 7.
\end{solution}

\begin{solution}{exo:b1:balanced-trees:3}
The tree: 5 with children 3 and 8; 3 with children 2 and 4; 1 the left child of 2. \omterm{def:b1:balanced-trees:avl}{Balance factors}: 1, 4 and 8 have 0;
2 has $-1 - 0 = -1$; 3 has $0 - 1 = -1$; 5 has $0 - 2 = -2$. It is not an \omterm{def:b1:balanced-trees:avl}{AVL tree}. The path from 5 to the deepest key
goes left twice (LL): a \omterm{def:b1:balanced-trees:rotation}{right rotation} at 5 gives root 3, children 2 (with 1) and 5 (with 4 and 8), all \omterm{def:b1:balanced-trees:avl}{balance factors}
in $\{-1, 0\}$.
\end{solution}

\begin{solution}{exo:b1:balanced-trees:4}
S: a black root. E: red left child of S. A: red left child of E, two reds in a row: \omterm{def:b1:balanced-trees:rotation}{right rotation} at S gives E with red
children A and S; the colour flip makes A and S black and E red, and the root is made black. R: larger than E, smaller
than S: red left child of S, which is allowed. Final tree: E (black) with children A (black) and S (black), and R (red)
the left child of S.
\end{solution}

\begin{solution}{exo:b1:balanced-trees:5}
Mia has two children; the shortcut returns its left child Eli, and the right \omterm{def:b1:induction-and-termination:rooted-tree}{subtree}, Tom with Pia and Zoe, is no longer
reachable: the tree keeps 3 keys of 6, and the three nodes of the lost \omterm{def:b1:induction-and-termination:rooted-tree}{subtree} are leaked, since nothing points to them
any more. The two-child case needs the \omterm{def:b1:balanced-trees:ordered}{successor}, as in \cref{lst:b1:balanced-trees:delete}.
\end{solution}

\begin{solution}{exo:b1:balanced-trees:6}
The \omterm{def:b1:balanced-trees:rotation}{rotation} at Eli makes Bea the root of that \omterm{def:b1:induction-and-termination:rooted-tree}{subtree}, with left child Abe and right child Eli, and empties Eli's left
link (Bea had no right child). But the call ignores the returned root: Mia still points to Eli, so Bea and Abe are lost
and the \omterm{def:b1:balanced-trees:traversal}{in-order traversal} prints Eli, Ivy, Mia and Tom. The repair stores the new root in the parent's link:
\code{t->left = rotate_right(t->left);}, as the AVL code does.
\end{solution}

\begin{solution}{exo:b1:balanced-trees:7}
\omcode{code/bachelor-1/17-balanced-trees/exercises/range_count.c}{6}{11}{Counting the keys in a range}
A node below \code{lo} has its whole left \omterm{def:b1:induction-and-termination:rooted-tree}{subtree} below \code{lo}; a node above \code{hi} its whole right \omterm{def:b1:induction-and-termination:rooted-tree}{subtree} above
\code{hi}. So the visited nodes are the ones in the range, plus at most two paths, along the boundaries: $O(h + r)$. The
test compares it with a count over a boolean \omterm{def:b1:pointers-and-arrays:array}{array} during 20\,000 random insertions and deletions.
\end{solution}

\begin{solution}{exo:b1:balanced-trees:8}
The tree with root 10, left child 5 and 20 the right child of 5 passes the local check (5 < 10, 20 > 5), but 20 lies in
the left \omterm{def:b1:induction-and-termination:rooted-tree}{subtree} of 10. The correct check passes an interval down:
\omcode{code/bachelor-1/17-balanced-trees/exercises/is_bst.c}{13}{20}{Checking with bounds}
A left child must lie between the lower bound of its parent and the parent's key, a right child between the parent's key
and its upper bound. The test confirms that the local check accepts the counterexample and the correct one rejects it.
\end{solution}

\begin{solution}{exo:b1:balanced-trees:9}
The program (\code{exercises/catalan.py}) finds 1, 1, 2, 5, 14 and 42 shapes for $n = 0, \dots, 5$, as does the
recurrence. A shape with $n \geq 1$ nodes is a root with a left shape of $i$ nodes and a right shape of $n - 1 - i$, for
$i = 0, \dots, n - 1$, independently: $C_n = \sum_i C_i C_{n-1-i}$. Every shape arises from some insertion order (insert
the keys in \omterm{def:b1:balanced-trees:traversal}{preorder} of the shape, labelled in order), so the program finds them all.
\end{solution}

\begin{solution}{exo:b1:balanced-trees:10}
By \omterm{def:b1:induction-and-termination:structural}{structural induction}, a \omterm{def:b1:induction-and-termination:rooted-tree}{subtree} whose root $x$ has \omterm{def:b1:balanced-trees:rb}{black height} $b$ (black nodes on a path from $x$ down, $x$
included if black) has at least $2^b - 1$ nodes. The empty \omterm{def:b1:induction-and-termination:rooted-tree}{subtree} has $b = 0$ and 0 nodes. Each child of $x$ has \omterm{def:b1:balanced-trees:rb}{black
height} $b - 1$ if $x$ is black and $b$ if $x$ is red, so at least $b - 1$: the \omterm{def:b1:induction-and-termination:rooted-tree}{subtree} has at least $1 + 2(2^{b-1} - 1) = 2^b - 1$ nodes. On a longest path, of $h + 1$ nodes from the root,
no two consecutive nodes are red and the root is black, so at least $\lceil (h + 1)/2 \rceil$ are black: $b \geq (h +
1)/2$. Hence $n \geq 2^{(h + 1)/2} - 1$, that is $h + 1 \leq 2 \log_2 (n + 1)$. The test checks it after every insertion.
\end{solution}

\begin{solution}{exo:b1:balanced-trees:11}
Let $z$ be the lowest node on the path that gets \omterm{def:b1:balanced-trees:avl}{balance factor} $-2$ after inserting into its left \omterm{def:b1:induction-and-termination:rooted-tree}{subtree} (the mirror
case is the same), and $h + 2$ its new height: before the insertion its left \omterm{def:b1:induction-and-termination:rooted-tree}{subtree} had height $h$, its right \omterm{def:b1:induction-and-termination:rooted-tree}{subtree}
$h - 1$, and $z$ height $h + 1$. LL case: the left child $y$ now has height $h + 1$, its left \omterm{def:b1:induction-and-termination:rooted-tree}{subtree} $h$, its right $h -
1$; after a \omterm{def:b1:balanced-trees:rotation}{right rotation} at $z$, $z$ has \omterm{def:b1:induction-and-termination:rooted-tree}{subtrees} of height $h - 1$ and $h - 1$, height $h$, and $y$ has height $h + 1$.
LR case: the right \omterm{def:b1:induction-and-termination:rooted-tree}{subtree} of $y$, rooted at $x$, has height $h$ and its \omterm{def:b1:induction-and-termination:rooted-tree}{subtrees} $h - 1$ and $h - 2$ in some order (both empty when $x$ is the new key, $h = 0$);
after the double \omterm{def:b1:balanced-trees:rotation}{rotation}, $x$ is the root with children $y$ and $z$ of height $h$ each: height $h + 1$. Either way the
\omterm{def:b1:induction-and-termination:rooted-tree}{subtree} has its height from before the insertion, so no ancestor's height or \omterm{def:b1:balanced-trees:avl}{balance factor} changes, and no other
\omterm{def:b1:balanced-trees:rotation}{rotation} happens. The test checks at most one \omterm{def:b1:balanced-trees:rotation}{rotation} event per insertion on 30 sequences.
\end{solution}

\begin{solution}{exo:b1:balanced-trees:12}
\omcode{code/bachelor-1/17-balanced-trees/exercises/from_sorted.c}{5}{12}{The middle key as the root}
Each key becomes one node, in constant time: $O(n)$. The two halves have $\lfloor (n - 1)/2 \rfloor$ and $\lceil (n -
1)/2 \rceil$ keys. By strong induction the height is $H(n) = \lfloor \log_2 n \rfloor$: $H(1) = 0$, and $H(n) = 1 +
H(\lceil (n-1)/2 \rceil) = 1 + \lfloor \log_2 \lfloor n/2 \rfloor \rfloor = \lfloor \log_2 n \rfloor$, since $\lceil
(n - 1)/2 \rceil = \lfloor n/2 \rfloor$. The test checks the height for every $n \leq 400$.
\end{solution}

\begin{solution}{pb:b1:balanced-trees:1}
\textbf{1.} Each new title is larger than all the previous ones: the search goes right at every node, and the new node
becomes the right child of the deepest one. The tree is a path of $n$ nodes.

\textbf{2.} The title at depth $d$ costs $d + 1$ comparisons: $(1 + 2 + \dots + 80)/80 = 40.5$.

\textbf{3.} The $i$-th insertion compares with the $i - 1$ titles already there: $0 + 1 + \dots + 79 = 3160$.

\textbf{4.} 11 against an average of 11.39 at $n = 80$, and close all along: the catalogue numbers are unrelated to the
alphabetical order of the titles, and behave like a random order.

\textbf{5.} $\lfloor \log_2 80 \rfloor = 6$: catalogue order 11, alphabetical 79, random 11.39 on average, AVL and
red-black 7.

\textbf{6.} Use a balanced tree, AVL or red-black; or rebuild a tree of minimal height from the sorted list each year
(\cref{exo:b1:balanced-trees:12}).

\textbf{7.} The new key is always the rightmost: the path from any unbalanced node to it goes right twice, the RR case,
repaired by a single left \omterm{def:b1:balanced-trees:rotation}{rotation}.

\textbf{8.} $22 + 11 = 33$, since an insertion makes at most one \omterm{def:b1:balanced-trees:rotation}{rotation} event. An insertion costs $O(\log n)$: the
descent and the way back up, $O(h)$ with $h \leq 1.44 \log_2 (n + 2)$, and a \omterm{def:b1:balanced-trees:rotation}{rotation} in $O(1)$.

\textbf{9.} 2 goes to the right of 1, and 3 has \omterm{def:b1:balanced-trees:avl}{balance factor} $-2$ with its left child 1 at $+1$: case LR. A left
\omterm{def:b1:balanced-trees:rotation}{rotation} at 1 makes 2 the left child of 3 with left child 1 (case LL), and a \omterm{def:b1:balanced-trees:rotation}{right rotation} at 3 gives root 2, children
1 and 3.

\textbf{10.} The \omterm{def:b1:balanced-trees:avl}{balance factors} are computed from the heights, which change only on the path of the insertion, and the
height of a node is $1 + \max$ of its children's stored heights: $O(1)$ per node.

\textbf{11.} No: $N(8) = 88 > 80$, so an \omterm{def:b1:balanced-trees:avl}{AVL tree} with 80 nodes has height at most 7. This one is as tall as possible.

\textbf{12.} After a deletion, the \omterm{def:b1:balanced-trees:rotation}{rotation} can leave the \omterm{def:b1:induction-and-termination:rooted-tree}{subtree} one level lower than before, which changes the \omterm{def:b1:balanced-trees:avl}{balance
factor} of its parent; the imbalance can then move up level by level.

\textbf{13.} $h + 1 = 8 \leq 2 \log_2 81 \approx 12.7$, and $2^5 - 1 = 31 \leq 80$.

\textbf{14.} The red nodes are F and I. The 2-3 tree has root D, children B and (F, H), and leaves A, C, E, G and
(I, J).

\textbf{15.} The 2-3 tree has $b$ levels and all its leaves at the same depth. With only one-key nodes it holds $1 + 2 +
\dots + 2^{b-1} = 2^b - 1$ keys; with only two-key nodes, level $i$ has $3^i$ nodes of 2 keys: $2(1 + 3 + \dots +
3^{b-1}) = 3^b - 1$.

\textbf{16.} $2^b - 1 \leq 80$ gives $b \leq 6$; $3^b - 1 \geq 80$ gives $b \geq 4$: $b$ is 4, 5 or 6. The tree has 5.

\textbf{17.} A red node does not change any \omterm{def:b1:balanced-trees:rb}{black height}; only the rule on red nodes can be broken, and the local repairs
fix it on the way up. A black node would break the equal \omterm{def:b1:balanced-trees:rb}{black heights}, which no local repair can fix.

\textbf{18.} \omterm{def:b1:balanced-trees:avl}{AVL trees} are lower ($1.44 \log_2 n$ against $2 \log_2 n$), so searches are slightly faster; an insertion
makes at most one \omterm{def:b1:balanced-trees:rotation}{rotation} event in both, but an AVL deletion may rotate at every level, a red-black deletion at most a
constant number of times. With many deletions, the \omterm{def:b1:balanced-trees:rb}{red-black tree}.

\textbf{19.} A tree of height $h \geq 2$ has a \omterm{def:b1:induction-and-termination:rooted-tree}{subtree} of height $h - 1$; the other has height $h - 1$ or $h - 2$, and the
fewest nodes come with $h - 2$; each \omterm{def:b1:induction-and-termination:rooted-tree}{subtree} must itself be a smallest \omterm{def:b1:balanced-trees:avl}{AVL tree} of its height, or a smaller one could
replace it. So $N(h) = 1 + N(h - 1) + N(h - 2)$. $N(0) = 1$ (a root) and $N(1) = 2$ (a root and a child).

\textbf{20.} $N(0) = 1 = F_3 - 1$, $N(1) = 2 = F_4 - 1$, and $N(h) = (F_{h+2} - 1) + (F_{h+1} - 1) + 1 = F_{h+3} - 1$.

\textbf{21.} $F_1 = 1 \geq \varphi^{-1}$ and $F_2 = 1 \geq \varphi^0$; then $F_{k+2} = F_{k+1} + F_k \geq \varphi^{k-1} +
\varphi^{k-2} = \varphi^{k-2}(\varphi + 1) = \varphi^k$.

\textbf{22.} $n \geq N(h) = F_{h+3} - 1 \geq \varphi^{h+1} - 1$, so $h + 1 \leq \log_\varphi (n + 1)$.

\textbf{23.} $\log_\varphi 81 = \ln 81 / \ln \varphi \approx 4.394/0.4812 \approx 9.13$, so $h \leq 8.13$, that is $h \leq
8$. The exact largest height is 7 ($N(7) = 54 \leq 80 < N(8) = 88$), and the measured \omterm{def:b1:balanced-trees:avl}{AVL tree} reaches it.

\textbf{24.} \emph{Minimal \omterm{def:b1:balanced-trees:avl}{AVL trees} are Fibonacci trees}: the smallest \omterm{def:b1:balanced-trees:avl}{AVL tree} of height $h$ is a root with smallest
\omterm{def:b1:induction-and-termination:rooted-tree}{subtrees} of heights $h - 1$ and $h - 2$, and it has $N(h) = F_{h+3} - 1$ nodes (questions 19 and 20; checked up to
$h = 20$, 28\,656 nodes). Hence every \omterm{def:b1:balanced-trees:avl}{AVL tree} with $n$ nodes has height at most $\log_\varphi (n + 1) - 1 < 1.45 \log_2
(n + 2)$: for the 80 titles, at most 8 by the bound and at most 7 exactly, the height measured in catalogue order.
\end{solution}
